\chapter{Cadre UE Chef d'oeuvre, Objectifs et Perimetre}

\section{Cadre pedagogique}
Ce document est redige dans le cadre de l'UE Chef d'oeuvre (IAFA, 2025--2026).
Il correspond au \textbf{Livrable 2}: document ecrit d'environ 20 pages,
centre sur les experimentations, evaluations et realisations effectuees.

\begin{table}[H]
  \centering
  \caption{Livrables et calendrier (consigne UE)}
  \begin{tabular}{p{2.1cm}p{8.4cm}p{3.2cm}}
    \toprule
    Livrable & Description & Date \\
    \midrule
    L1 & Premier rapport (etat de l'art, methodes, premieres experimentations, gestion de projet) & 31 dec. 2025 \\
    L2 & Second rapport (experimentations, evaluations, realisations) & 17 mars 2026 \\
    L3 & Soutenance orale (20 min + 15 min questions) & 19/20 mars 2026 \\
    L4 & Demonstration + depot code GitHub formation (non obligatoire alternants) & selon consigne \\
    \bottomrule
  \end{tabular}
\end{table}

\begin{table}[H]
  \centering
  \caption{Ponderation annoncee pour la note S10}
  \begin{tabular}{lr}
    \toprule
    Composante & Poids \\
    \midrule
    Note de travail encadrant & 30\% \\
    Livrable 2 & 30\% \\
    Livrable 4 & 10\% \\
    Soutenance & 30\% \\
    \bottomrule
  \end{tabular}
\end{table}

\section{Objectifs techniques du projet}
Le projet repond a l'enonce ``Generation automatisee de narratives analogues via LLM''
avec les exigences suivantes:
\begin{itemize}[leftmargin=2em]
  \item exploitation du corps des articles Wikipedia (pas uniquement les infobox),
  \item gestion des textes longs par chunking,
  \item prise en compte du multilingue FR/EN,
  \item apprentissage d'un sentence-encoder via AutoTrain,
  \item interface Spaces avec generation iterative conditionnee par un score de similarite.
\end{itemize}

\section{Perimetre retenu}
Le corpus cible est centre sur des personnalites (principalement acteurs/actrices).
Le pipeline retenu pour les resultats finaux s'appuie sur des paires
\texttt{en <-> fr\_en} (traduction FR->EN), avec negatives difficiles (hard negatives).

\section{Hors perimetre de cette version}
\begin{itemize}[leftmargin=2em]
  \item extension systematique aux categories ``films'' et ``livres'',
  \item annotation humaine exhaustive de qualite factuelle,
  \item optimisation cout/latence en production.
\end{itemize}
